\subsection{漏服之后怎么办：口服避孕药的差错处理}

要说清"漏服"，开头必须先分清药的类型——不同类型药的容错窗口差别极大，用错了规则，等于没有规则。

\textbf{一、先确认你吃的是哪一类}

\begin{itemize}
\item \textbf{复方口服避孕药（含雌激素 + 孕激素）}：常见为 21 片装（服 21 天、停 7 天）、24$+$4 片装或 28 片连续装。它的容错窗口是 \textbf{12 小时}——部分含新型孕激素的配方为 24 小时，以自己药盒内的说明书为准；
\item \textbf{单纯孕激素避孕药（微丸，mini-pill）}：只含孕激素、不含雌激素。它的窗口通常只有 \textbf{3 小时}（少数新一代配方为 12 小时）——\textbf{"12 小时规则"对它完全不适用}。这类药对服药时刻的要求更严格，属于"差一小时都可能掉保护"的一类。
\end{itemize}

所以本节的第一个提醒是：\emph{先看药盒，再套用下面的规则}。

\textbf{二、复方口服避孕药的漏服规则}

\begin{itemize}
\item \textbf{延迟不足 12 小时}：立即补服，下一片仍按原定时间服用，\textbf{无需额外措施}；
\item \textbf{延迟 12--24 小时（漏服 1 片）}：立即补服（可能一次吃两片），并在\textbf{此后 7 天加用屏障避孕}（避孕套）；
\item \textbf{漏服 2 片及以上}：补服最近漏掉的一片，其余漏服的药片丢弃——\textbf{不要一次补吃多片}；此后 7 天加用屏障避孕。另需按漏服发生的周次分别处理：
  \begin{itemize}
  \item 若漏服发生在\textbf{第一周}，且在此前 5 天内有过无保护性行为——考虑紧急避孕（见本书紧急避孕章节）；
  \item 若漏服发生在\textbf{第三周}：跳过停药期，直接开始服下一盒（或按说明书"补足后不停药"的方案），以免停药期与漏服叠加，形成一个足够长的排卵窗口；
  \end{itemize}
\item \textbf{补服后出现点滴出血}：这是激素波动引起的撤退性出血，不是"药失效"的证据，也不是停药的理由。
\end{itemize}

\textbf{三、单纯孕激素避孕药的漏服规则}

\begin{itemize}
\item \textbf{延迟超过 3 小时}：立即补服，此后 \textbf{48 小时}加用屏障避孕；
\item 若在漏服前后有过无保护性行为，同样需要考虑紧急避孕；
\item 由于窗口极短，这类药更适合作息规律、能固定时刻服药的人。若经常错过，应与医生讨论换成贴片、阴道环、皮下埋植或宫内节育器。
\end{itemize}

\textbf{四、呕吐与腹泻：看不见的漏服}

\begin{itemize}
\item \textbf{服药后 3--4 小时内呕吐}：药片可能尚未吸收，按"漏服 1 片"处理，补服 1 片；
\item \textbf{严重腹泻}（大量水样便、持续一天以上）：吸收可能受影响，腹泻期间及症状结束后 7 天建议加用屏障避孕；
\item 反复呕吐、腹泻且无法补服时，可咨询医生是否需要临时改用其他方式。
\end{itemize}

\textbf{五、把"漏服"变成少见事件}

\begin{itemize}
\item \textbf{绑定日常动作}：把药放在牙刷旁、早餐旁等每天都会在同一时间经过的位置，比单纯"记得吃"有效得多；
\item \textbf{双重提醒}：设两个提醒——服药时刻一个、15 分钟后的补漏提醒一个；
\item \textbf{诚实评估自己}：如果半年内已漏服多次，那么"我适合每天吃药"这个前提本身就不成立。换成长效方法（贴片、阴道环、皮下埋植、宫内节育器）不是失败，而是把方法调到适合自己的位置（参见"如何选择合适的避孕方法"）；
\item \textbf{不必自责}：漏服是所有每日服药方案共有的问题，不是纪律问题。把它当成一个设计问题来解决，比要求自己"下次一定记得"更可靠。
\end{itemize}

\textbf{六、拿不准的时候怎么办}

\begin{itemize}
\item 原则很简单：\textbf{拿不准，就当成漏了}。补服 1 片、加用屏障 7 天，这两步成本极低；赌一次的代价（无准备的怀孕，或紧急避孕）要高得多；
\item \textbf{不要自行加量}：除说明书规定的补服规则外，不要一次服用更多片来"补偿"；
\item \textbf{不要自行停药}：中途停药会引起撤退性出血并失去保护。如想更换方法，应与医生商量过渡方案；
\item 若已连续漏服两天以上、又不确定是否已进入排卵窗口，最稳妥的做法是：立即补服、7 天屏障避孕，并在本次月经推迟超过一周时做妊娠测试。
\end{itemize}

\textit{【编者注】}以上是通用规则。不同配方的容错窗口可能不同，具体以自己药盒内的说明书与医生意见为准。若同时在服用可能降低避孕药效果的药物（部分抗癫痫药、利福平等），漏服的影响会被进一步放大，建议提前与医生确认是否需要加用屏障。

